\chapter{Contexte du stage}
\label{chap:contexte}

\section*{Introduction}
\addcontentsline{toc}{section}{Introduction}

Ce chapitre situe le stage dans son cadre. Il présente l'organisme d'accueil, l'objectif
poursuivi, et surtout le périmètre volontairement borné du projet --- car dans un stage de huit
semaines, ce qu'on décide de \emph{ne pas} faire compte autant que ce qu'on entreprend. Il décrit
enfin l'organisation du travail et l'instrument qui l'a rendue traçable : un journal de bord tenu
quotidiennement.

\section{Organisme d'accueil}

Pulsaride Solutions est une société française du domaine \anglais{deeptech}, dont le siège est
situé rue Henri Dunant à Pessac, en Gironde.

\begin{table}[htbp]
  \centering
  \caption{Fiche signalétique de l'organisme d'accueil}
  \label{tab:entreprise}
  \small
  \begin{tabularx}{\textwidth}{@{}>{\raggedright\arraybackslash}p{3.4cm} R@{}}
    \toprule
    \sffamily\bfseries\color{navy}Élément &
    \sffamily\bfseries\color{navy}Description \\
    \midrule
    Raison sociale   & Pulsaride Solutions \textsc{sarl} \\
    Siège            & Pessac, Gironde, France \\
    Domaine          & Cybersécurité et conformité réglementaire \\
    Spécialité       & Attestation cryptographique du comportement des systèmes logiciels et des agents d'intelligence artificielle \\
    Marché visé      & Institutions financières européennes soumises aux réglementations DORA, NIS2 et au règlement européen sur l'intelligence artificielle \\
    Encadrant        & M. Lyazid \textsc{Salihi}, fondateur et dirigeant \\
    \bottomrule
  \end{tabularx}
\end{table}

L'entreprise développe une plateforme de défense opérant au niveau du noyau du système
d'exploitation, qui capture les appels système, les signe cryptographiquement et produit des
certificats vérifiables \emph{hors ligne} --- c'est-à-dire sans qu'aucune donnée ne soit confiée à
un service tiers pour être contrôlée.

Ce détail n'est pas anecdotique pour le présent stage : il explique le sujet. La proposition de
l'entreprise tient en une idée --- fournir une \textbf{preuve vérifiable} plutôt qu'une assurance
sur parole, et la fournir sans dépendance à une infrastructure distante. Le sujet confié
transpose exactement cette idée au domaine juridique : un assistant qui \emph{cite ses sources}
pour qu'on puisse les vérifier, et qui s'exécute \emph{localement} pour que la question ne quitte
pas la machine. Les deux exigences structurantes de ce rapport --- vérifiabilité et souveraineté
--- sont donc celles de l'entreprise, appliquées à un autre objet.

Le stage s'est ainsi inscrit dans l'axe «~intelligence artificielle appliquée~» de l'entreprise,
sur un sujet articulant traitement automatique du langage et souveraineté des données.
L'encadrement a été assuré par M. Lyazid \textsc{Salihi}.

\section{Objectif et missions}

L'objectif est de concevoir et d'évaluer un assistant de recherche juridique souverain fondé sur
l'architecture RAG. Le citoyen pose une question en langage naturel ; le système recherche le
passage pertinent dans le corpus juridique, puis formule une réponse claire \textbf{en citant
systématiquement la source officielle}.

Le principe de fiabilité est central, et il a une formulation opérationnelle précise : le système
s'appuie sur les textes réels et \textbf{indique honnêtement lorsqu'il ne dispose pas de
l'information}, plutôt que de produire une réponse vraisemblable.

\begin{encadre}[title=Exemple d'interaction visée]
\textbf{Citoyen} \quad «~J'ai droit à combien de jours de congé de maternité~?~»

\medskip
\textbf{Assistant} \quad «~Selon l'article~152 du Code du travail (Loi~65-99), la salariée
bénéficie d'un congé de maternité de quatorze semaines. Le contrat de travail est suspendu, non
rompu, pendant cette période.

\smallskip
\textit{Source : Code du travail, article~152.}

\smallskip
Réponse fournie à titre informatif ; pour une situation particulière, consultez un
professionnel du droit.~»
\end{encadre}

Le stage s'articule autour de quatre missions.

\begin{enumerate}
  \item \textbf{Constitution et préparation du corpus juridique.} Collecter, nettoyer et découper
        le corpus par article. La qualité de ce découpage conditionne directement la qualité des
        réponses : un article tronqué produit une réponse tronquée, sans qu'aucun signal ne le
        révèle en aval.
  \item \textbf{Moteur de recherche sémantique et génération.} Indexer le corpus, retrouver les
        passages pertinents à partir d'une question, puis générer une réponse claire et sourcée,
        en garantissant l'ancrage dans les textes réels et la gestion explicite du «~je ne sais
        pas~».
  \item \textbf{Évaluation de la fiabilité.} Construire un jeu de questions-réponses de référence
        et mesurer la précision du système. L'analyse honnête des limites fait partie intégrante
        du livrable, au même titre que les résultats favorables.
  \item \textbf{Extension multilingue} \textit{(ambition complémentaire)}. Explorer une couche de
        normalisation linguistique permettant de comprendre une question posée en arabe ou en
        darija, tout en conservant l'ancrage des réponses dans le texte officiel d'origine.
\end{enumerate}

\section{Périmètre et non-objectifs}

Pour garantir la faisabilité en huit semaines, le sujet a été volontairement borné : un
\textbf{périmètre juridique unique} --- le droit du travail --- et un fonctionnement
\textbf{monolingue solide} en français. L'extension multilingue et l'élargissement à d'autres
codes constituent des ambitions complémentaires, non des engagements.

Le cadrage négatif est ici aussi important que le cadrage positif. Le système n'est explicitement
pas :

\begin{itemize}
  \item \textbf{un outil de conseil juridique personnalisé.} Il relève de la \emph{restitution
        d'information juridique} ; chaque réponse porte un avertissement en ce sens ;
  \item \textbf{un système qui répond «~au mieux~».} Face à une question sans réponse dans le
        corpus, \textbf{l'abstention est le comportement correct}, pas un échec ;
  \item \textbf{un produit fini.} Le livrable est le moteur RAG, son évaluation et le dépôt
        reproductible qui les accompagne.
\end{itemize}

Quatre interdictions ont été posées dès la conception et tenues jusqu'au terme du stage : pas de
\anglais{fine-tuning} du modèle, pas de mémoire conversationnelle multi-tours, pas d'autre code
juridique que le Code du travail, et aucune intégration d'interface de programmation externe dans
le cœur du système. Chacune écarte une dérive dont le coût aurait été payé en fin de stage, au
moment où il n'aurait plus été possible de revenir en arrière.

\section{Organisation du travail}

Le travail a été planifié en huit phases hebdomadaires, chacune assortie d'un livrable identifié à
l'avance. Le découpage suivait une progression délibérée : rendre le corpus fiable avant de
construire quoi que ce soit dessus, obtenir une chaîne complète avant de l'optimiser, et mesurer
avant de conclure.

L'exécution réelle s'est nettement écartée de ce découpage, et c'est l'écart qui est instructif.
Le journal de bord permet de le dater précisément : les trois premières semaines, du 2 au
20~juillet, ont été absorbées par le \emph{seul} corpus --- extraction, découpage, reconstruction
de la hiérarchie, puis correction des familles de défauts décrites au
chapitre~\ref{chap:realisation}. Le moteur lui-même, du premier pipeline RAG bout-en-bout jusqu'à
l'évaluation complète, a été mené en deux jours, les 24 et 25~juillet.

\begin{constat}
Ce déséquilibre --- trois semaines de données contre deux jours de moteur --- n'est pas un
accident de planification. Il est la manifestation directe de la thèse défendue dans ce rapport :
une fois le corpus fiable, la boucle RAG se construit vite. C'est la \textbf{donnée} qui coûte,
pas le code qui l'exploite.
\end{constat}

Un \textbf{journal de bord} a été tenu quotidiennement, consignant chaque décision, chaque
résultat et chaque difficulté --- y compris les impasses et les corrections de corrections. Ce
journal a servi de matière première à la rédaction du présent rapport, ce qui explique que les
échecs y soient rapportés avec autant de précision que les succès : ils y avaient été écrits sur
le moment, avant qu'on en connaisse l'issue.

\section*{Conclusion}
\addcontentsline{toc}{section}{Conclusion}

Le projet répond à un besoin concret --- rendre le droit du travail interrogeable en langage
naturel --- sous une contrainte structurante, la souveraineté des données, et avec une exigence
centrale, la fiabilité. Le périmètre a été borné dès le premier jour, et les quatre interdictions
posées à ce moment-là n'ont pas été levées. Le chapitre suivant expose les fondements techniques
sur lesquels le système repose, en commençant par la raison pour laquelle un modèle de langage
seul ne suffit pas.
